\documentclass[12pt,a4paper]{article}
\usepackage[UTF8]{ctex}
\usepackage{geometry}
\usepackage{booktabs}
\usepackage{multirow}
\usepackage{amsmath}
\usepackage{caption}
\usepackage{array}
\geometry{left=2.5cm,right=2.5cm,top=2.5cm,bottom=2.5cm}

\title{CPR-GMRES求解器单进程性能对比测试报告}
\author{}
\date{测试日期：2026年6月9日}

\begin{document}
\maketitle

\section{实验目的}

对比以下三种CPR-GMRES求解器在单进程（NP=1）下的求解性能，分析混合精度策略对求解效率和收敛性的影响：
\begin{itemize}
  \item \textbf{原始双精度求解器}（\texttt{test\_bsr\_cprgmres}）：全双精度 CPR-GMRES，收敛容差 $10^{-4}$
  \item \textbf{混合精度 GMRES 求解器}（\texttt{test\_bsr\_cprgmres\_hybrid}）：外层双精度 GMRES，内层单精度 CPR 预处理，收敛容差 $10^{-6}$
  \item \textbf{迭代精化求解器}（\texttt{test\_bsr\_cprgmres\_mixed}）：双精度外残差校正 + 单精度内层 GMRES 求解的 IR 方案，收敛容差 $10^{-6}$
\end{itemize}

\section{实验环境}

\begin{itemize}
  \item 硬件平台：XeonGold5122
  \item 编译器：mpicc（GCC 6.1.0），开启 -O3 优化
  \item MPI环境：Intel MPI 2021.6.0，单进程运行
\end{itemize}

\section{测试问题}

\begin{itemize}
  \item 测试矩阵：SPE10 油藏模型，BSR格式，块大小 $2 \times 2$
  \item 标量维度：$2\,188\,844 \times 2\,188\,844$
  \item 非零块数：$7\,478\,310$
  \item 解耦方法：True-IMPES
  \item Krylov子空间维度：30
\end{itemize}

\section{三种cpr-gmres方法核心性能指标对比}

本实验在单进程（NP=1）下运行，针对同一油藏模拟线性系统，分别测试三种 CPR-GMRES 求解器的性能，stage2使用GS磨光2次。
\begin{table}[htbp]
\centering
\caption{三种求解器单进程性能对比}
\label{tab:compare}
\begin{tabular}{lrrr}
\toprule
指标 & 原始双精度 & Hybrid & Mixed (IR) \\
\midrule
\textbf{Setup 阶段} \\
~ Decoupling 时间 (s) & 0.1053 & 0.1011 & 0.0848 \\
~ CPR Setup 时间 (s)  & 1.9519 & 1.4817 & 1.4310 \\
~ GMRES Setup 时间 (s)& 0.0029 & ---    & 0.0028 \\
~ Setup 总时间 (s)    & 2.0602 & 1.5828 & 1.5185 \\
\midrule
\textbf{Solve 阶段} \\
~ Precond 总时间 (s)  & 27.29  & 18.05  & --- \\
~ Solve 时间 (s)      & 34.62  & 25.90  & 53.46 \\
\midrule
\textbf{总计} \\
~ Total 时间 (s)      & 36.68  & 27.38  & 54.98 \\
\midrule
\textbf{收敛信息} \\
~ 迭代次数            & 25     & 30     & 19 (IR) \\
~ 收敛容差            & $10^{-4}$ & $10^{-6}$ & $10^{-6}$ \\
~ 最终相对残量        & $9.13\times10^{-5}$ & $2.76\times10^{-5}$ & $6.15\times10^{-5}$ \\
\bottomrule
\end{tabular}
\end{table}

  从上表可以看出，Hybrid 方案仅通过将 CPR 预处理单精度化，即以轻微的迭代次数增加换取了显著的每步加速，整体有效；而 Mixed IR 将整个内层求解器下沉至单精度，虽然降低了单次内迭代开销，但受限于内迭代次数过多，总代价超过原始双精度方案。

\newpage
\section{逐迭代步收敛对比}

表~\ref{tab:gmres_residual} 给出了三种 CPR-GMRES 求解器在单进程求解过程中，外层迭代步数与相对残量范数的变化历程。

\begin{table}[htbp]
\centering
\caption{逐迭代步相对残量对比（Original \& Hybrid）}
\label{tab:gmres_residual}
\begin{tabular}{crrr}
\toprule
迭代步 & 原始双精度 & Hybrid \\
\midrule
 1 & $2.3164\times10^{-1}$ & $2.0430\times10^{-1}$ & $1.6765\times10^{-1}$ \\
 5 & $1.2934\times10^{-2}$ & $1.3293\times10^{-2}$ & $2.1748\times10^{-2}$ \\
10 & $3.2809\times10^{-3}$ & $3.3352\times10^{-3}$ & $2.3449\times10^{-3}$ \\
15 & $1.1715\times10^{-3}$ & $1.1647\times10^{-3}$ & $3.2562\times10^{-4}$ \\
20 & $3.5912\times10^{-4}$ & $3.5883\times10^{-4}$ & $6.1462\times10^{-5}$ \\
25 & $9.1281\times10^{-5}$ & $9.2186\times10^{-5}$ & 上面为第19步 \\
30 & ---                  & $2.7649\times10^{-5}$ \\
\bottomrule
\end{tabular}
\end{table}

Hybrid 仅将 CPR 预处理下沉至单精度，外层 GMRES 仍以双精度维持收敛方向，故残差历史与原始方案高度一致，仅需少量额外迭代即可达标，同时单次迭代成本显著降低，整体收益为正。而 Mixed IR 将整个求解器置于单精度内层，外迭代虽在双精度校正下收敛迅速，但内层精度不足导致每次外迭代需多次内迭代才能满足内收敛要求，总计算量过大，抵消了单精度加速的收益。

\newpage

\section{mixed策略内外迭代明细}

表~\ref{tab:ir_inner_detail} 展示了混合精度迭代精化（Mixed IR）求解器在单进程求解过程中，前 3 次 IR 外迭代的内层 GMRES 逐步残量明细。
\begin{table}[htpb]
    \centering
    \caption{Mixed (IR) 前3步内外迭代逐步残量明细}
    \label{tab:ir_inner_detail}
    \begin{tabular}{ccll}
    \toprule
    IR 外迭代 & 内层 GMRES 步 & 残量范数 & 相对残量（对内层右端项） \\
    \midrule
    \multirow{5}{*}{0}
     & 1 & $2.7128\times10^{+0}$ & $2.9995\times10^{-1}$ \\
     & 2 & $2.2572\times10^{+0}$ & $2.4957\times10^{-1}$ \\
     & 3 & $1.9863\times10^{+0}$ & $2.1962\times10^{-1}$ \\
     & 4 & $1.7299\times10^{+0}$ & $1.9127\times10^{-1}$ \\
     & 5 & $1.5146\times10^{+0}$ & $1.6746\times10^{-1}$ \\
    \cmidrule(lr){1-4}
     & \multicolumn{3}{l}{$\rightarrow$ IR 外残量：$1.5163\times10^{+0}$（相对原始右端项：$1.6765\times10^{-1}$）} \\
    \midrule
    \multirow{5}{*}{1}
     & 1 & $1.3417\times10^{+0}$ & $8.8494\times10^{-1}$ \\
     & 2 & $1.2039\times10^{+0}$ & $7.9400\times10^{-1}$ \\
     & 3 & $1.0657\times10^{+0}$ & $7.0290\times10^{-1}$ \\
     & 4 & $9.1582\times10^{-1}$ & $6.0402\times10^{-1}$ \\
     & 5 & $8.0779\times10^{-1}$ & $5.3277\times10^{-1}$ \\
    \cmidrule(lr){1-4}
     & \multicolumn{3}{l}{$\rightarrow$ IR 外残量：$8.0930\times10^{-1}$（相对原始右端项：$8.9483\times10^{-2}$）} \\
    \midrule
    \multirow{5}{*}{2}
     & 1 & $7.2440\times10^{-1}$ & $8.9514\times10^{-1}$ \\
     & 2 & $6.5271\times10^{-1}$ & $8.0656\times10^{-1}$ \\
     & 3 & $5.9554\times10^{-1}$ & $7.3592\times10^{-1}$ \\
     & 4 & $5.4717\times10^{-1}$ & $6.7614\times10^{-1}$ \\
     & 5 & $4.9559\times10^{-1}$ & $6.1240\times10^{-1}$ \\
    \cmidrule(lr){1-4}
     & \multicolumn{3}{l}{$\rightarrow$ IR 外残量：$4.9570\times10^{-1}$（相对原始右端项：$5.4809\times10^{-2}$）} \\
    \bottomrule
    \end{tabular}
    \end{table}
    
    内层不充分的收敛是 Mixed IR 在当前配置下低效的根源。若欲使 IR 方案有效，必须显著增强内层求解能力（如增加内迭代次数或引入更有效的预条件），或结合多右端技术摊薄内层开销，否则难以在单右端问题中体现性能优势。

    \newpage

    \section{混合精度求解器hybrid并行扩展性（NP=1 vs NP=4）}

    表~\ref{tab:hybrid_scaling} 展示了混合精度 GMRES 求解器（Hybrid）在单进程（NP=1）与四进程（NP=4）配置下的性能对比。测试算例与前文一致，收敛容差均为 $10^{-6}$。主要计时指标定义同表~\ref{tab:compare}。加速比按 NP=1 时间除以 NP=4 时间计算，反映从单进程到四进程的强可扩展性。

\begin{table}[htbp]
\centering
\caption{Hybrid 求解器 NP=1 与 NP=4 性能对比}
\label{tab:hybrid_scaling}
\begin{tabular}{lrrr}
\toprule
指标 & NP=1 & NP=4 & 加速比 \\
\midrule
\textbf{Setup 阶段} \\
~ Decoupling 时间 (s) & 0.1011 & 0.0783 & 1.29$\times$ \\
~ CPR Setup 时间 (s)  & 1.4817 & 0.8963 & 1.65$\times$ \\
~ Setup 总时间 (s)    & 1.5828 & 0.9746 & 1.62$\times$ \\
\midrule
\textbf{Solve 阶段} \\
~ Precond 总时间 (s)  & 18.05  & 6.42   & 2.81$\times$ \\
~ Solve 时间 (s)      & 25.90  & 8.20   & 3.16$\times$ \\
\midrule
\textbf{总计} \\
~ Total 时间 (s)      & 27.38  & 9.10   & 3.01$\times$ \\
\midrule
\textbf{收敛信息} \\
~ 迭代次数            & 30     & 30     & --- \\
~ 最终相对残量        & $2.76\times10^{-5}$ & $2.60\times10^{-5}$ & --- \\
\bottomrule
\end{tabular}
\end{table}

将 NP=4 下的 Hybrid 总时间（9.10 s）与 NP=1 下的原始双精度总时间（36.68 s，见表~\ref{tab:compare}）相比，混合精度策略与四进程并行共同带来了约 4.03 倍的整体加速。两者优势可有效叠加，进一步验证了 Hybrid 方案作为并行基准的潜力。

\newpage

\section{混合精度求解器hybrid并行扩展性（NP=1 vs NP=4）}

表~\ref{tab:mixed_scaling} 对比了混合精度迭代精化（Mixed IR）求解器在单进程（NP=1）与四进程（NP=4）下的性能表现。两者均采用双精度外残差校正与单精度内层 GMRES 求解的 IR 框架，收敛容差设为 $10^{-6}$，每次外迭代固定执行 5 次单精度内迭代。

\begin{table}[htbp]
    \centering
    \caption{Mixed (IR) 求解器 NP=1 与 NP=4 性能对比}
    \label{tab:mixed_scaling}
    \begin{tabular}{lrrr}
    \toprule
    指标 & NP=1 & NP=4 & 加速比 \\
    \midrule
    \textbf{Setup 阶段} \\
    ~ Decoupling 时间 (s) & 0.0848 & 0.0802 & 1.06$\times$ \\
    ~ CPR Setup 时间 (s)  & 1.4310 & 1.5299 & 0.94$\times$ \\
    ~ Setup 总时间 (s)    & 1.5185 & 1.6131 & 0.94$\times$ \\
    \midrule
    \textbf{Solve 阶段} \\
    ~ Precond 总时间 (s)  & ---    & ---    & --- \\
    ~ Solve 时间 (s)      & 53.46  & 82.25  & 0.65$\times$ \\
    \midrule
    \textbf{总计} \\
    ~ Total 时间 (s)      & 54.98  & 83.86  & 0.66$\times$ \\
    \midrule
    \textbf{收敛信息} \\
    ~ 迭代次数            & 19 (IR) & 50 (IR) & --- \\
    ~ 最终相对残量        & $6.15\times10^{-5}$ & $7.30\times10^{-2}$ & --- \\
    \bottomrule
    \end{tabular}
    \end{table}

    并行计算引入的浮点运算非确定性（如求和次序变化、通信舍入误差）使内层单精度求解提供的校正量精度进一步劣化，外残差方程无法被足够精确地近似满足，导致外迭代失去收敛驱动力。这一数值失稳将 IR 方案的“可能缓慢收敛”变为“完全不可靠”。
    
\clearpage
\section{附录：可调参数参考}

以下列出当前CPR-GMRES求解器中所有可调参数，按模块分类。标注"*"者为当前测试中实际使用的值。

\subsection{命令行参数（所有求解器共用）}

\begin{table}[htbp]
\centering
\caption{命令行参数}
\label{tab:params_cmdline}
\begin{tabular}{llll}
\toprule
参数 & 位置 & 默认值 & 说明 \\
\midrule
\texttt{matrix\_file} & argv[1] & --- & 矩阵文件路径 \\
\texttt{matrix\_type} & argv[2] & --- & 矩阵类型：0=CSR, 1=BSR \\
\texttt{binary} & argv[3] & 0 & 文件格式：0=文本, 1=二进制 \\
\texttt{rhs\_file} & argv[4] & NULL & 右端项文件路径（可选） \\
\bottomrule
\end{tabular}
\end{table}

\subsection{GMRES 求解器参数}

\subsubsection*{原始双精度求解器（\texttt{test\_bsr\_cprgmres.c}）}

\begin{table}[htbp]
\centering
\caption{原始双精度求解器 GMRES 参数（代码中硬编码）}
\label{tab:params_orig_gmres}
\begin{tabular}{llll}
\toprule
参数名 & 类型 & 当前值 * & 说明 \\
\midrule
\texttt{k\_dim} & int & 30 & Krylov子空间维度 \\
\texttt{max\_iter} & int & 100 & 最大GMRES迭代次数 \\
\texttt{tol} & double & 1e-4 & 相对残差收敛容差 \\
\texttt{print\_level} & int & 1 & 打印级别（0=不打印, 1=基本, 2=详细） \\
\texttt{is\_check\_restarted} & int & 1 & 重启后检查收敛（0=关闭, 1=开启） \\
\bottomrule
\end{tabular}
\end{table}

\subsubsection*{混合精度 GMRES（\texttt{test\_bsr\_cprgmres\_hybrid.c}）}

\begin{table}[htbp]
\centering
\caption{Hybrid 求解器 GMRES 参数}
\label{tab:params_hybrid_gmres}
\begin{tabular}{llll}
\toprule
参数名 & 类型 & 当前值 * & 说明 \\
\midrule
\texttt{k\_dim} & int & 30 & Krylov子空间维度 \\
\texttt{max\_iterations} & int & 30 & 最大GMRES迭代次数 \\
\texttt{tolerance} & double & 1e-6 & 相对残差收敛容差 \\
\texttt{print\_level} & int & 1 & 打印级别 \\
\bottomrule
\end{tabular}
\end{table}

\subsubsection*{IR 混合精度（\texttt{test\_bsr\_cprgmres\_mixed.c}）}

\begin{table}[htbp]
\centering
\caption{IR 求解器参数}
\label{tab:params_ir}
\begin{tabular}{llll}
\toprule
参数名 & 类型 & 当前值 * & 说明 \\
\midrule
\multicolumn{4}{l}{\textbf{外层 IR 参数}} \\
\texttt{max\_ir\_iterations} & int & 50 & 最大IR外迭代次数 \\
\texttt{ir\_tolerance} & double & 1e-4 & 外层IR收敛容差 \\
\multicolumn{4}{l}{\textbf{内层单精度 GMRES 参数}} \\
\texttt{k\_dim} & int & 30 & Krylov子空间维度 \\
\texttt{max\_iter} & int & 5 & 内层GMRES最大迭代次数 \\
\texttt{tol} & float & 1e-2 & 内层GMRES收敛容差 \\
\texttt{print\_level} & int & 1 & 打印级别 \\
\bottomrule
\end{tabular}
\end{table}

\subsection{CPR 预条件器参数}

所有三个求解器均通过 \texttt{JXF\_CPRSetParameter()} 和 \texttt{JXF\_CPRSetRealParameter()} 设置CPR参数。

\begin{table}[htbp]
\centering
\caption{CPR 预条件器参数}
\label{tab:params_cpr}
\begin{tabular}{llll}
\toprule
参数名 & 类型 & 当前值 * & 说明 \\
\midrule
\texttt{pressure\_index} & int & 0 & 压力变量在块中的索引 \\
\texttt{block\_size} & int & 2 & 矩阵块大小 \\
\texttt{stage1\_maxit} & int & 1 & 阶段1(AMG压力求解)最大迭代次数 \\
\texttt{stage2\_maxit} & int & 2/3* & 阶段2(块磨光)迭代次数 \\
\texttt{stage1\_solver\_type} & int & 1 & 阶段1求解器类型：1=AMG \\
\texttt{stage2\_solver\_type} & int & 2/3* & 阶段2类型：2=BGS, 3=BILU-GMRES \\
\texttt{print\_level} & int & 1/2* & CPR打印级别 \\
\texttt{threshold} & float & 1e-15 & 提取压力矩阵的阈值 \\
\bottomrule
\end{tabular}
\end{table}

\subsection{PAMG（并行代数多重网格）参数}

CPR的阶段1使用PAMG求解压力子系统。以下参数通过 \texttt{JXF\_PAMGSet*()} 系列函数设置，若未显式设置则使用默认值。

\begin{table}[htbp]
\centering
\caption{PAMG 核心参数}
\label{tab:params_pamg_core}
\begin{tabular}{llll}
\toprule
参数函数 & 类型 & 默认值 & 说明 \\
\midrule
\texttt{SetMaxLevels} & int & 25 & 最大网格层数 \\
\texttt{SetStrongThreshold} & float & 0.25 & 强连接阈值（越小则粗化越慢） \\
\texttt{SetMaxRowSum} & float & 0.9 & 最大行和阈值（用于强度矩阵） \\
\texttt{SetTruncFactor} & float & 0.0 & 插值截断因子 \\
\texttt{SetPMaxElmts} & int & 5 & 插值每行最大元素数 \\
\texttt{SetCoarsenType} & int & 6 & 粗化策略：0=CLJP, 6=Falgout, 8=PMIS, 10=HMIS \\
\texttt{SetInterpType} & int & 0 & 插值类型：0=经典, 1=直接, 3=多通, 6=ExtPI \\
\texttt{SetMeasureType} & int & 0 & 度量方式：0=局部, 1=全局 \\
\texttt{SetCycleType} & int & 1 & 循环类型：1=V-cycle, 2=W-cycle \\
\texttt{SetNumSweeps} & int & 1 & 每层松弛扫描次数 \\
\texttt{SetRelaxType} & int & 3 & 松弛方法：0=Jacobi, 3=hybrid GS(前向), \\
  &   &   & 4=hybrid GS(后向), 6=hybrid SSOR \\
\texttt{SetRelaxOrder} & int & 1 & 松弛顺序：0=自然序, 1=CF-松弛 \\
\texttt{SetRelaxWt} & float & 1.0 & 松弛权重（0=自动计算） \\
\texttt{SetOuterWt} & float & 1.0 & SOR/SSOR外部松弛权重 \\
\texttt{SetTol} & float & 1e-7 & PAMG作为求解器时的容差 \\
\texttt{SetCoarsestSolverID} & int & 9 & 最粗层求解器：9=高斯消元 \\
\texttt{SetCoarseThreshold} & int & 100 & 最粗层允许的最大网格点数 \\
\texttt{SetCoarseRatio} & float & 0.75 & 粗化比例阈值 \\
\texttt{SetNumFunctions} & int & 1 & PDE函数数（多物理场问题） \\
\texttt{SetNodal} & int & 0 & 节点型粗化（0=标量, 1=节点） \\
\bottomrule
\end{tabular}
\end{table}

\begin{table}[htbp]
\centering
\caption{PAMG 高级平滑器参数}
\label{tab:params_pamg_smooth}
\begin{tabular}{llll}
\toprule
参数函数 & 类型 & 默认值 & 说明 \\
\midrule
\texttt{SetSmoothType} & int & 6 & 复杂平滑器类型：6=Schwarz, 7=Pilut, \\
  &   &   & 8=ParaSails, 9=Euclid \\
\texttt{SetSmoothNumLevels} & int & 0 & 使用复杂平滑器的层数（0=不使用） \\
\texttt{SetSmoothNumSweeps} & int & 1 & 复杂平滑器扫描次数 \\
\texttt{SetVariant} & int & 0 & Schwarz变体（0=乘性, 1=加性） \\
\texttt{SetOverlap} & int & 1 & Schwarz重叠层数 \\
\texttt{SetDomainType} & int & 1 & 子域类型 \\
\texttt{SetSchwarzRlxWeight} & float & 1.0 & Schwarz松弛权重 \\
\bottomrule
\end{tabular}
\end{table}

\begin{table}[htbp]
\centering
\caption{PAMG 聚合式多重网格参数}
\label{tab:params_pamg_agg}
\begin{tabular}{llll}
\toprule
参数函数 & 类型 & 默认值 & 说明 \\
\midrule
\texttt{SetAggNumLevels} & int & 0 & 聚合层数（0=不使用聚合） \\
\texttt{SetAggInterpType} & int & 默认 & 聚合插值类型 \\
\texttt{SetAggPMaxElmts} & int & 默认 & 聚合插值每行最大元素 \\
\texttt{SetAggTruncFactor} & float & 默认 & 聚合插值截断因子 \\
\bottomrule
\end{tabular}
\end{table}

\subsection{解耦（Decoupling）参数}

\begin{table}[htbp]
\centering
\caption{解耦参数}
\label{tab:params_decoup}
\begin{tabular}{llll}
\toprule
参数名 & 类型 & 当前值 * & 说明 \\
\midrule
\texttt{decoup\_type} & enum & TIMPES(5) & 解耦方法枚举： \\
  &   &   & 0=None, 1=ABF, 2=ANL, 3=SEM,  \\
  &   &   & 4=QI, 5=TIMPES, 6=TIMPES2 \\
\texttt{is\_thermal} & int & 0 & 是否为热力模型（影响解耦算法） \\
\bottomrule
\end{tabular}
\end{table}

\subsection{ILU/BILU 预条件器参数}

阶段2使用BILU-GMRES时的可调参数（原始双精度求解器中使用）：

\begin{table}[htbp]
\centering
\caption{ILU/BILU 参数}
\label{tab:params_ilu}
\begin{tabular}{llll}
\toprule
参数函数 & 类型 & 默认值 & 说明 \\
\midrule
\texttt{SetMaxIter} & int & 默认 & ILU求解最大迭代次数 \\
\texttt{SetDropTol} & float & 0.0 & ILU分解丢弃阈值 \\
\texttt{SetNxy} & int[2] & --- & 网格维度（nx, ny） \\
\texttt{SetNpxy} & int[2] & --- & 进程划分（npx, npy） \\
\texttt{SetNumEquns} & int & 1 & 方程数（多物理场） \\
\bottomrule
\end{tabular}
\end{table}

\subsection{GMRES 收敛判据参数（高级）}

GMRES求解器中可选的收敛判据参数，见 \texttt{jxf\_krylov.h} 中 \texttt{jxf\_GMRESData} 结构体：

\begin{table}[htbp]
\centering
\caption{GMRES 高级收敛参数}
\label{tab:params_gmres_adv}
\begin{tabular}{llll}
\toprule
参数函数 & 类型 & 默认值 & 说明 \\
\midrule
\texttt{SetMinIter} & int & 0 & 最小迭代次数 \\
\texttt{SetAbsoluteTol} & float & 0.0 & 绝对残差容差 \\
\texttt{SetConvergenceFactorTol} & float & 1.0 & 收敛因子阈值 \\
\texttt{SetStopCrit} & int & 0 & 停止准则 \\
\texttt{SetConvCriteria} & int & 0 & 收敛判据：0=相对残差范数 \\
\texttt{SetConvFacThreshold} & float & 1.0 & 收敛因子终止阈值 \\
\texttt{SetRelChange} & int & 0 & 是否检查相对变化 \\
\texttt{SetLogging} & int & 1 & 日志级别（0=关闭, 1=基本, 2=详细） \\
\bottomrule
\end{tabular}
\end{table}

\subsection{运行脚本配置参数（\texttt{run\_test.sh} / \texttt{cpr\_test.sh}）}

\begin{table}[htbp]
\centering
\caption{运行脚本配置参数}
\label{tab:params_script}
\begin{tabular}{llll}
\toprule
参数名 & 默认值 & 说明 \\
\midrule
\texttt{NP\_LIST} & "1 2 4 8 16 32" & MPI进程数列表 \\
\texttt{METHOD} & "BILU-GMRES" & 方法标识（用于日志文件名） \\
\texttt{DECOUP\_TYPE} & "5" (TIMPES) & 解耦类型（设为 "all" 遍历全部） \\
\texttt{MATRIX\_FILE} & 见脚本 & 矩阵文件路径 \\
\texttt{MATRIX\_TYPE} & "1" (BSR) & 矩阵类型 \\
\texttt{BINARY} & "1" (二进制) & 文件格式 \\
\texttt{RHS\_FILE} & 见脚本 & 右端项文件 \\
\bottomrule
\end{tabular}
\end{table}

\subsection{调试与输出参数}

\begin{table}[htbp]
\centering
\caption{调试与输出参数}
\label{tab:params_debug}
\begin{tabular}{llll}
\toprule
参数函数/变量 & 类型 & 默认值 & 说明 \\
\midrule
\texttt{print\_level} & int & 1 & 全局打印级别 \\
\texttt{logging} & int & 1 & 日志记录级别 \\
\texttt{SetLogging} & --- & 1 & PAMG日志（0=关闭, 1=基本） \\
\texttt{SetPrintLevel} & --- & 1 & PAMG打印（0=无, 1=setup, 2=solve, 3=全部） \\
\texttt{SetDebugFlag} & --- & 0 & 调试标记 \\
\texttt{SetPrintFileName} & char* & NULL & 输出文件名 \\
\texttt{SetPrintCoarseMatrix} & int & 0 & 是否打印粗网格矩阵 \\
\texttt{log\_file\_name} & char* & NULL & GMRES日志文件名 \\
\bottomrule
\end{tabular}
\end{table}

以上参数覆盖了从命令行输入、求解器算法、预条件器配置到输出调试的全链路可调选项。实际使用中，可通过修改对应源文件中的参数值来调整求解行为。

\end{document}
